%% ============================================================
%% 分片 01：原文第 1 节「Introduction」（原文 glshort.tex 起于第 313 行）
%% 由 tools/harvmac2tex.py 从 arXiv 源码（hep-th/0604151v3）自动生成。
%% 只译英文散文；公式、编号、\label、\cite 一律不改。
%% ============================================================

\section{Introduction}
\label{sec:intro}

 The Langlands program for number fields \cite{ref1} unifies many classical
and contemporary results in number theory and is a vast area of
research. 

It has an analog for curves over a finite field, which has also
been the subject of much celebrated work \cite{ref2,ref3}. In addition, a geometric version of the Langlands
program for curves has been much developed 

 
\cite{ref4,ref5,ref6,ref7,ref8}, both for curves over a field of
characteristic $p$ and for ordinary complex Riemann surfaces. For
a survey that is relatively readable for physicists, with numerous
references, see \cite{ref9}.

Our focus in the present paper is on the geometric Langlands
program for complex Riemann surfaces.  We aim to show how this
program can be understood as a chapter in quantum field theory. No
prior familiarity with the Langlands program is assumed; instead,
we assume a familiarity with subjects such as supersymmetric gauge
theories, electric-magnetic duality, sigma-models, mirror
symmetry, branes, and topological field theory. The theme of the
paper is to show that when these familiar physical ingredients are
applied to just the right problem, the geometric Langlands program
arises naturally. Seemingly esoteric notions such as Hecke
eigensheaves, ${\cal D}$-modules, and so on, appear spontaneously
in the physics, with new insights about their properties.

The first hints of a connection between the Langlands program and
quantum field theory came from the work of Goddard, Nuyts, and
Olive (GNO), who showed in 1976 \cite{ref10} that in gauge theories, though electric
charge takes values in the weight lattice of the gauge group,
magnetic charge takes values in the weight lattice of a dual
group. Magnetic charges for general compact Lie groups had been
first analyzed by Englert and Windey \cite{ref11}. The
GNO analysis motivated the Montonen-Olive electric-magnetic
duality conjecture \cite{ref12} according to which a specific gauge theory based
on a given gauge group is equivalent to a similar theory with the
coupling constant inverted and the gauge group replaced by its
dual.

For a gauge group $G$, the GNO dual group is actually the same as
the Langlands dual group $^L\negthinspace G$, which plays an
important role in formulating the Langlands conjectures. (For some
examples of the correspondence between $G$ and $^L\negthinspace
G$, see the table.) This was observed by Atiyah, who suggested to
the second author at the end of 1977 that the Langlands program is
related to quantum field theory and recommended the two papers
\cite{ref10,ref12}. There resulted a further development
\cite{ref13} in which it was understood that Montonen-Olive duality is
more natural in supersymmetric gauge theory.  It was later
understood that ${\EUN}=4$ supersymmetry (i.e., the maximal
possible supersymmetry in four dimensions) is the right context
for this duality \cite{ref14}, and that the
$\Bbb{Z}_2$ duality originally proposed has a natural extension to
$SL(2,\Bbb{Z})$ 

\cite{ref15,ref16} when the theta angle of the gauge theory is
included.\par\medskip \begin{table}[htbp]
\centering
\begin{tabular}{ll}
\toprule
$G$ & ${}^L G$ \\
\midrule
$U(N)$ & $U(N)$ \\
$SU(N)$ & $PSU(N)=SU(N)/\Bbb{Z}_N$ \\
$Spin(2n)$ & $SO(2n)/\Bbb{Z}_2$ \\
$Sp(n)$ & $SO(2n+1)$ \\
$Spin(2n+1)$ & $Sp(n)/\Bbb{Z}_2$ \\
$G_2$ & $G_2$ \\
$E_8$ & $E_8$ \\
\bottomrule
\end{tabular}
\caption{Examples of the correspondence between a Lie group $G$ and its Langlands or GNO dual ${}^L\negthinspace G$.（原文 Table~1）}
\label{tab:gno-dual}
\end{table}

 

 In the early 1990's, extensions of
Montonen-Olive duality to string theory were conjectured
\cite{ref17}. Subsequently, Montonen-Olive duality and its
generalizations were studied from many new points of view and were
recognized as a crucial, though still mysterious, ingredient in
understanding field theory and string theory.  These developments
were far too extensive to be reviewed here, but one observation of
that period, though a sideline at the time, is a starting point
for the present paper. Compactification of ${\EUN}=4$ super
Yang-Mills theory from four dimensions to two dimensions was
studied \cite{ref18,ref19} and was shown to lead at low energies to
a two-dimensional supersymmetric sigma-model in which the target
space is a hyper-Kahler manifold that is Hitchin's moduli space
${\EUM}_H$ of stable Higgs bundles
 \cite{ref20}.
 Electric-magnetic duality in four dimensions
reduces in two dimensions to $T$-duality of the sigma-model.  This
particular $T$-duality was subsequently used mathematically
\cite{ref21} to show (for $SU(N)$) that
the Hodge numbers of the Higgs bundle moduli space of a gauge
group $G$ are equal to those of $^L\negthinspace G$. The geometry
underlying this $T$-duality was investigated in \cite{ref22} and
subsequently in \cite{ref23} for any
semi-simple Lie group $G$.

Other clues about the relation of the geometric Langlands program
to quantum field theory have come from relatively recent
mathematical work. The approach of Beilinson and Drinfeld to the
geometric Langlands program is based on quantization of $\MH$, as
the title of their paper implies \cite{ref5}. The $T$-duality of
${\EUM}_H$, understood mathematically as a Fourier-Mukai
transform, has been interpreted as a sort of semiclassical
approximation to the geometric Langlands program. This point of
view underlies the paper \cite{ref24}. We understand that there have also been
important unpublished contributions by other mathematicians,
including Donagi and Pantev. The second author learned of this
interpretation of the $T$-duality of $\MH$ from a lecture by D.
Ben-Zvi at a conference on the geometric Langlands program, held
at the IAS in the spring of 2004. This was an extremely strong
hint that it must be possible to understand the geometric
Langlands program using four-dimensional electric-magnetic duality
(which leads to this particular $T$-duality) and branes (the
natural quantum field theory setting for interpreting $T$-duality
as a Fourier-Mukai transform).  This hint was the starting point
for the present paper.

To summarize this paper in the briefest possible terms, we will
develop six main ideas.   The first is that from a certain twisted
version of ${\EUN}=4$ supersymmetric Yang-Mills theory in four
dimensions, one can construct a family of four-dimensional
topological field theories in four dimensions. After reviewing
some of the background in section 第~\ref{sec:review}~节, we explain this
construction in section 第~\ref{sec:tqftfour}~节. The twisting procedure is
formally just analogous to the construction by which Donaldson
theory can be obtained \cite{ref25}\ from ${\EUN}=2$ super Yang-Mills
theory. The second main idea, developed in sections
第~\ref{sec:compacthitchin}~节\ and 第~\ref{sec:tqfttwo}~节, is that, extending the insights of
\cite{ref18,ref19}, compactification on a Riemann surface $C$ gives
in two dimensions a family of topological sigma-models, with
target ${\EUM}_H$, which are ``generalized $B$-models.'' Moreover,
for a special value of the parameter, four-dimensional $S$-duality
acts as two-dimensional mirror symmetry.  The third main idea,
developed in section 第~\ref{sec:loopop}~节, is that Wilson and 't Hooft line
operators are topological operators that act on the branes of the
two-dimensional sigma-model in a natural fashion. Here we consider
an operator that maps one brane to another (or roughly speaking,
one theory to another), not the more familiar sort of operator in
Hilbert space that maps one state to another. A brane that is
mapped by the Wilson or 't Hooft operators to, roughly speaking, a
multiple of itself is what we call an electric or magnetic
eigenbrane. $S$-duality will automatically exchange the electric
and magnetic eigenbranes. The fourth main idea, explained in
section 第~\ref{sec:electricbranes}~节, is that, in the right context, electric
eigenbranes are in natural correspondence with homomorphisms of
$\pi_1(C)$ to the complexification of the Langlands dual group
$^L\negthinspace G_{\Bbb{C}}$. The fifth main idea, developed in
sections 第~\ref{sec:thooftheckeop}~节\ and 第~\ref{sec:bogeqheck}~节, is that 't Hooft
operators correspond naturally to geometric Hecke operators
similar to those of the geometric Langlands program but acting on
Higgs bundles instead of ordinary $G$-bundles. It takes one more
important idea, developed in section 第~\ref{sec:abranes}~节, to make contact
with the usual formulation of the geometric Langlands
correspondence. We show that, because of the existence of a
canonical coisotropic brane on ${\EUM}_H$, the magnetic
eigenbranes in our sense are naturally associated to ${\cal
D}$-modules (modules for the sheaf of differential operators) on
the moduli space ${\EUM}(G,C)$ of holomorphic $G$-bundles on $C$.

Sections 7, 8, 9, and 10 of this paper and some expository
portions of other sections are primarily adapted from a
forthcoming book \cite{ref26} which will also contain
some additional results. This material is included here to make
the paper more comprehensible.

 
 

 

    

 
 
 One obvious gap in our analysis is
that we consider only the unramified case of the geometric
Langlands correspondence.  We expect that it is possible to apply
somewhat similar ideas to the ramified case.  A second major gap
is that we do not shed light on the utility of two-dimensional
conformal field theory for the geometric Langlands program
\cite{ref5,ref9,ref27,ref28,ref29,ref30}. The last of these
references applies conformal field theory to the ramified case.
Hopefully it will prove possible to deduce the conformal field
theory approach from the gauge theory approach of this paper. In
fact, there is an analogy even at a naive level \cite{ref31}\ between
conformal field theory and the theory of automorphic
representations, which is the basis of the Langlands program.
Finally, though we have nothing to contribute about this here, an
additional clue about the relation of the Langlands program with
physics presumably comes from the diverse ways that automorphic
forms enter string theory. For a tiny sampling of this, see
\cite{ref32,ref33,ref34,ref35,ref36,ref37,ref38,ref39,ref40}.

A.K. would like to thank D. Arinkin, R. Bezrukavnikov, and D.
Orlov for useful conversations. In particular, Orlov's
explanations in 2002 about the abelian case of the geometric
Langlands program partially motivated the paper \cite{ref41}, which will enter our story in section 11.

E.W. would like to thank the many mathematicians who over the
years have explained matters relevant to the Langlands program,
including A. Beilinson, P. Deligne,  V. Drinfeld,  and K. Vilonen,
and especially M. F. Atiyah,   D. Ben-Zvi, R. Donagi, E. Frenkel,
and D. Kazhdan, and most recently M. Goresky and R. MacPherson. In
addition, among others,  T. Hausel, N. Hitchin, M. Hopkins, P.
Kronheimer, L. Jeffrey,  J. Morgan, G. Moore, D. Morrison, N.
Nekrasov, M. Thaddeus, C. Vafa, and E. J. Weinberg clarified some
points relevant to the present paper, and many of the physicists
at the IAS, including S. Hellerman, K. Intriligator, J. Maldacena,
N. Seiberg, and J. Walcher, made helpful comments.
